\makeatletter
\def\input@path{{./}{../}{../../}{../../../}{../../../../}{preamble/}{../preamble/}{../../preamble/}{../../../preamble/}{../../../../preamble/}}
\makeatother

\input{preamble_G10_T1_C2C3.tex}

\title{Grade 10 Term 1 Lesson 1 Practice Activity 2\\Revenue, Costs, Profit, and Payoff Tables}
\author{}
\date{}

\begin{document}
	\maketitle

	\setcounter{tocdepth}{1}
	\setcounter{secnumdepth}{2}
	\phantomsection
	\label{toc}
	\tableofcontents

	\section*{Evaluated criteria}
	\begin{enumerate}[label=\textbf{C\arabic*}]
		\item Characterises decision-making problems under uncertainty by identifying decision alternatives, states of nature, consequences, and the economic or financial context of the decision.
		\item[\textbf{C3}] Calculates profit as revenue minus costs across states of nature and organises the resulting payoffs in payoff tables.
	\end{enumerate}

	\section{Introduction}
	\label{sec:introduction}

	\SectionProblemsTableOfContents{
		\SectionProblemLink{sec:problem1-introductory-general-c1-c3}{Introductory General Problem (C1--C3)} \\
		\SectionProblemLink{sec:problem2-advanced-general-c1-c3}{Advanced General Problem (C1--C3)}
	}

	\subsection{Introductory General Problem (C1--C3)}
	\label{sec:problem1-introductory-general-c1-c3}

	A school enterprise must commit to either alternative \(A_1\) or alternative \(A_2\) before learning whether demand will be \(S_1\) or \(S_2\). Its accounting team used sales forecasts for each choice under each demand outcome to estimate total revenue, and used purchasing and operating budgets to estimate the corresponding total cost. All amounts and the resulting profit are measured in thousand USD.

	If the enterprise selects \(A_1\), forecast sales produce total revenue of \(20\) under demand state \(S_1\) and \(13\) under \(S_2\); fulfilling those sales is budgeted to cost \(9\) and \(8\), respectively. If it selects \(A_2\), its different sales plan is expected to bring in \(18\) under \(S_1\) and \(15\) under \(S_2\), while the resources required by that plan cost \(10\) and \(9\), respectively. The enterprise must calculate the profit for every possible alternative--state pair before constructing its payoff table.

	\textbf{Questions/Tasks.}
	\begin{enumerate}
		\CriterionTask
		{subsec:problem1-introductory-general-c1}{1}
		{Interpret the decision elements.}

		\CriterionTask
		{subsec:problem1-introductory-general-c3}{3}
		{Build the payoff table.}
	\end{enumerate}

	\subsubsection*{C1 - Interpreting the decision.}
	\phantomsection
	\label{subsec:problem1-introductory-general-c1}

	\DecisionElementsTable
	{\(A_1\); \(A_2\)}
	{\(S_1\); \(S_2\)}
	{Demand revealed after the enterprise choice}
	{Profit in thousand USD for each alternative--state pair}

	\secInicio[sec:problem1-introductory-general-c1-c3][sec:introduction]

	\subsubsection*{C3 - Calculating revenue, cost, and payoff.}
	\phantomsection
	\label{subsec:problem1-introductory-general-c3}

	\begin{center}
		\small
		\begin{tabular}{lrrr}
			\toprule
			Cell
			& \(R(A_i,S_j)\)
			& \(C(A_i,S_j)\)
			& \(P(A_i,S_j)=R(A_i,S_j)-C(A_i,S_j)\) \\
			\midrule
			\((A_1,S_1)\) & 20 & 9  & \(20-9=11\) \\
			\((A_1,S_2)\) & 13 & 8  & \(13-8=5\) \\
			\((A_2,S_1)\) & 18 & 10 & \(18-10=8\) \\
			\((A_2,S_2)\) & 15 & 9  & \(15-9=6\) \\
			\bottomrule
		\end{tabular}
	\end{center}

	\begin{center}
		\begin{tabular}{lrr}
			\toprule
			Alternative & \(S_1\) & \(S_2\) \\
			\midrule
			\(A_1\) & 11 & 5 \\
			\(A_2\) & 8  & 6 \\
			\bottomrule
		\end{tabular}
	\end{center}

	\secInicio[sec:problem1-introductory-general-c1-c3][sec:introduction]

	\subsection{Advanced General Problem (C1--C3)}
	\label{sec:problem2-advanced-general-c1-c3}

	A community business must choose operating plan \(A_1\) or \(A_2\) before demand is revealed as \(S_1\) or \(S_2\). Its demand study indicates that plan \(A_1\) will complete \(2\) units in state \(S_1\) and \(1\) in state \(S_2\), whereas plan \(A_2\) will complete \(3\) and \(2\) units, respectively. All revenue rates, cost rates, fixed amounts, and profits are measured in thousand USD, with each variable rate applying to every unit completed.

	The plans command different base selling rates because they offer different services: each unit under \(A_1\) earns an alternative-based revenue rate of \(10\), while each unit under \(A_2\) earns \(11\). Demand conditions add a market-based revenue rate of \(5\) per unit in \(S_1\) or \(3\) per unit in \(S_2\). The signed agreements also provide plan-based fixed revenue of \(7\) for \(A_1\) and \(9\) for \(A_2\), while community funding adds state-based fixed revenue of \(6\) in \(S_1\) and \(7\) in \(S_2\).

	The business's operating budget assigns an alternative-based cost of \(3\) per unit to \(A_1\) and \(4\) per unit to \(A_2\), reflecting the resources used by each plan. Demand-related handling adds a further state-based cost of \(1\) per unit in \(S_1\) and \(2\) per unit in \(S_2\). Preparing the plans also creates fixed costs of \(8\) for \(A_1\) and \(10\) for \(A_2\), and state-specific administration creates fixed costs of \(5\) in \(S_1\) and \(3\) in \(S_2\). The business must combine all four revenue components and all four cost components for each pair, then subtract total cost from total revenue.

	\textbf{Questions/Tasks.}
	\begin{enumerate}
		\CriterionTask
		{subsec:problem2-advanced-general-c1}{1}
		{Interpret the decision elements.}

		\CriterionTask
		{subsec:problem2-advanced-general-c3}{3}
		{Build the payoff table.}
	\end{enumerate}

	\subsubsection*{C1 - Interpreting the decision.}
	\phantomsection
	\label{subsec:problem2-advanced-general-c1}

	\DecisionElementsTable
	{\(A_1\); \(A_2\)}
	{\(S_1\); \(S_2\)}
	{Demand revealed after the business choice}
	{Profit in thousand USD from the complete revenue and cost models}

	For every alternative--state pair, revenue and cost are calculated independently:
	\[
	\begin{aligned}
		R(A_i,S_j)
		={}&
		q(A_i,S_j)
		\left[
		r(A_i)+r(S_j)
		\right]
		+
		f_r(A_i)
		+
		f_r(S_j),
		\\
		C(A_i,S_j)
		={}&
		q(A_i,S_j)
		\left[
		c(A_i)+c(S_j)
		\right]
		+
		f_c(A_i)
		+
		f_c(S_j),
		\\
		P(A_i,S_j)
		={}&
		R(A_i,S_j)-C(A_i,S_j).
	\end{aligned}
	\]

	\secInicio[sec:problem2-advanced-general-c1-c3][sec:introduction]

	\subsubsection*{C3 - Calculating revenue, cost, and payoff.}
	\phantomsection
	\label{subsec:problem2-advanced-general-c3}

	\begin{center}
		\begin{tabular}{llll}
			\toprule
			Cell
			& \(R(A_i,S_j)\)
			& \(C(A_i,S_j)\)
			& \(P(A_i,S_j)=R(A_i,S_j)-C(A_i,S_j)\) \\
			\midrule
			\((A_1,S_1)\)
			& \(2(10+5)+7+6=43\)
			& \(2(3+1)+8+5=21\)
			& \(43-21=22\) \\
			
			\((A_1,S_2)\)
			& \(1(10+3)+7+7=27\)
			& \(1(3+2)+8+3=16\)
			& \(27-16=11\) \\
			
			\((A_2,S_1)\)
			& \(3(11+5)+9+6=63\)
			& \(3(4+1)+10+5=30\)
			& \(63-30=33\) \\
			
			\((A_2,S_2)\)
			& \(2(11+3)+9+7=44\)
			& \(2(4+2)+10+3=25\)
			& \(44-25=19\) \\
			\bottomrule
		\end{tabular}
	\end{center}

	\begin{center}
		\begin{tabular}{lrr}
			\toprule
			Alternative & \(S_1\) & \(S_2\) \\
			\midrule
			\(A_1\) & 22 & 11 \\
			\(A_2\) & 33 & 19 \\
			\bottomrule
		\end{tabular}
	\end{center}

	\secInicio[sec:problem2-advanced-general-c1-c3][sec:introduction]

	\section{Direct Revenue and Cost Models}
	\label{sec:direct-revenue-cost-models}

	\SectionProblemsTableOfContents{
		\SectionProblemLink{sec:problem3-local-cafe-investment-c1-c3}{Local Café Investment (C1--C3)} \\
		\SectionProblemLink{sec:problem4-green-energy-production-mix-c1-c3}{Green Energy Production Mix (C1--C3)}
	}

	\subsection{Local Café Investment (C1--C3)}
	\label{sec:problem3-local-cafe-investment-c1-c3}

	A small café must decide whether to launch a line of artisan desserts or keep its current menu before learning whether daily customer traffic will be high or low. The owner prepared direct total-revenue forecasts from expected café sales and direct total-cost budgets covering ingredients, staffing, preparation, and other resources. Every amount and the resulting daily profit are in thousand USD.

	Launching the dessert line could generate revenue of \(58\) when traffic is high, when more customers buy the new products, and \(27\) when traffic is low. The matching operating budgets are \(16\) and \(19\); the low-traffic cost remains comparatively high because ingredients and preparation must be arranged before actual traffic is known. Keeping the current menu instead is forecast to produce revenue of \(49\) under high traffic and \(41\) under low traffic, with established-menu operating costs of \(21\) in either state. The owner must calculate profit for all four outcomes.

	\textbf{Questions/Tasks.}
	\begin{enumerate}
		\CriterionTask
		{subsec:problem3-local-cafe-investment-c1}{1}
		{Interpreting decision alternatives, events, consequences, and states.}

		\CriterionTask
		{subsec:problem3-local-cafe-investment-c3}{3}
		{Calculating profit and building the payoff table.}
	\end{enumerate}

	\subsubsection*{C1 - Interpreting decision alternatives, events, consequences, and states.}
	\phantomsection
	\label{subsec:problem3-local-cafe-investment-c1}

	\DecisionElementsTable
	{Launch dessert line; Keep current menu}
	{High traffic; Low traffic}
	{Actual foot traffic after the decision}
	{Profit in thousand USD for each alternative--state pair}

	\secInicio[sec:problem3-local-cafe-investment-c1-c3][sec:direct-revenue-cost-models]

	\subsubsection*{C3 - Calculating profit and building the payoff table.}
	\phantomsection
	\label{subsec:problem3-local-cafe-investment-c3}

	For every alternative--state pair, calculate
	\[
	P(A_i,S_j)=R(A_i,S_j)-C(A_i,S_j).
	\]

	\begin{center}
		\small
		\begin{tabular}{lrrr}
			\toprule
			Cell
			& \(R(A_i,S_j)\)
			& \(C(A_i,S_j)\)
			& \(P(A_i,S_j)=R(A_i,S_j)-C(A_i,S_j)\) \\
			\midrule
			\((A_1,S_1)\) & 58 & 16 & \(58-16=42\) \\
			\((A_1,S_2)\) & 27 & 19 & \(27-19=8\) \\
			\((A_2,S_1)\) & 49 & 21 & \(49-21=28\) \\
			\((A_2,S_2)\) & 41 & 21 & \(41-21=20\) \\
			\bottomrule
		\end{tabular}
	\end{center}

	The resulting payoff table is:

	\begin{center}
		\begin{tabular}{lrr}
			\toprule
			Alternative & High traffic & Low traffic \\
			\midrule
			Launch dessert line & 42 & 8 \\
			Keep current menu & 28 & 20 \\
			\bottomrule
		\end{tabular}
	\end{center}

	\secInicio[sec:problem3-local-cafe-investment-c1-c3][sec:direct-revenue-cost-models]

	\subsection{Green Energy Production Mix (C1--C3)}
	\label{sec:problem4-green-energy-production-mix-c1-c3}

	A renewable-energy company must build wind turbines, solar farms, or a balanced hybrid system before the year's weather is known. Annual electricity sales and the operating requirements of each generation system have been forecast directly for windy, sunny, and cloudy conditions. All revenue, cost, and profit figures are in thousand USD.

	Wind turbines are expected to earn revenue of \(132\) in windy weather, \(78\) in sunny weather, and \(62\) in cloudy weather because their electricity output changes with conditions; the corresponding operating and maintenance costs are \(42\), \(33\), and \(32\). Solar farms instead are forecast to earn \(77\), \(143\), and \(63\) in the same respective states, with costs of \(42\), \(48\), and \(38\), reflecting the resources needed at each output level.

	The hybrid system spreads production across both sources, producing forecast revenue of \(112\) when it is windy, \(107\) when it is sunny, and \(96\) when it is cloudy. Its corresponding operating costs are \(42\), \(42\), and \(36\). Management must subtract the appropriate direct cost from revenue for each strategy--weather combination.

	\textbf{Questions/Tasks.}
	\begin{enumerate}
		\CriterionTask
		{subsec:problem4-green-energy-production-mix-c1}{1}
		{Interpreting decision alternatives, events, consequences, and states.}

		\CriterionTask
		{subsec:problem4-green-energy-production-mix-c3}{3}
		{Calculating profit and building the payoff table.}
	\end{enumerate}

	\subsubsection*{C1 - Interpreting decision alternatives, events, consequences, and states.}
	\phantomsection
	\label{subsec:problem4-green-energy-production-mix-c1}

	\DecisionElementsTable
	{Wind turbines; Solar farms; Hybrid system}
	{Windy; Sunny; Cloudy}
	{Weather outcome affecting electricity generation}
	{Profit in thousand USD for each strategy--weather combination}

	\secInicio[sec:problem4-green-energy-production-mix-c1-c3][sec:direct-revenue-cost-models]

	\subsubsection*{C3 - Calculating profit and building the payoff table.}
	\phantomsection
	\label{subsec:problem4-green-energy-production-mix-c3}

	For every alternative--state pair, calculate
	\[
	P(A_i,S_j)=R(A_i,S_j)-C(A_i,S_j).
	\]

	\begin{center}
		\small
		\begin{tabular}{lrrr}
			\toprule
			Cell
			& \(R(A_i,S_j)\)
			& \(C(A_i,S_j)\)
			& \(P(A_i,S_j)=R(A_i,S_j)-C(A_i,S_j)\) \\
			\midrule
			\((A_1,S_1)\) & 132 & 42 & \(132-42=90\) \\
			\((A_1,S_2)\) & 78  & 33 & \(78-33=45\) \\
			\((A_1,S_3)\) & 62  & 32 & \(62-32=30\) \\
			\((A_2,S_1)\) & 77  & 42 & \(77-42=35\) \\
			\((A_2,S_2)\) & 143 & 48 & \(143-48=95\) \\
			\((A_2,S_3)\) & 63  & 38 & \(63-38=25\) \\
			\((A_3,S_1)\) & 112 & 42 & \(112-42=70\) \\
			\((A_3,S_2)\) & 107 & 42 & \(107-42=65\) \\
			\((A_3,S_3)\) & 96  & 36 & \(96-36=60\) \\
			\bottomrule
		\end{tabular}
	\end{center}

	The resulting payoff table is:

	\begin{center}
		\begin{tabular}{lrrr}
			\toprule
			Strategy & Windy & Sunny & Cloudy \\
			\midrule
			Wind turbines & 90 & 45 & 30 \\
			Solar farms & 35 & 95 & 25 \\
			Hybrid system & 70 & 65 & 60 \\
			\bottomrule
		\end{tabular}
	\end{center}

	\secInicio[sec:problem4-green-energy-production-mix-c1-c3][sec:direct-revenue-cost-models]

	\section{One-Component Models}
	\label{sec:one-component-models}

	\SectionProblemsTableOfContents{
		\SectionProblemLink{sec:problem5-custom-printing-product-line-c1-c3}{Custom Printing Product Line (C1--C3)} \\
		\SectionProblemLink{sec:problem6-event-equipment-fleet-c1-c3}{Event Equipment Fleet (C1--C3)} \\
		\SectionProblemLink{sec:problem7-community-arts-sponsorship-c1-c3}{Community Arts Sponsorship (C1--C3)} \\
		\SectionProblemLink{sec:problem8-distribution-contract-fulfilment-c1-c3}{Distribution Contract Fulfilment (C1--C3)}
	}

	\subsection{Custom Printing Product Line (C1--C3)}
	\label{sec:problem5-custom-printing-product-line-c1-c3}

	A custom-printing cooperative must choose to produce premium notebooks or canvas bags before it knows whether the market will bring strong demand with stable material prices or weak demand with higher material prices. Orders already under discussion indicate that premium notebooks will sell \(6\) production lots in the strong state and \(4\) in the weak state, while canvas bags will sell \(5\) and \(3\) lots, respectively.

	Because the two products have different selling prices, each notebook lot earns an alternative-based revenue rate of \(12\) thousand USD and each canvas-bag lot earns \(15\) thousand USD. Suppliers quote by market condition rather than product: stable prices make the cost \(5\) thousand USD per lot in the strong state, while higher prices raise it to \(8\) thousand USD per lot in the weak state. These are the only revenue and cost components, so the cooperative must apply both rates to the appropriate forecast lot quantity.

	\textbf{Questions/Tasks.}
	\begin{enumerate}
		\CriterionTask
		{subsec:problem5-custom-printing-product-line-c1}{1}
		{Interpreting decision alternatives, events, consequences, and states.}

		\CriterionTask
		{subsec:problem5-custom-printing-product-line-c3}{3}
		{Calculating profit and building the payoff table.}
	\end{enumerate}

	\subsubsection*{C1 - Interpreting decision alternatives, events, consequences, and states.}
	\phantomsection
	\label{subsec:problem5-custom-printing-product-line-c1}

	\DecisionElementsTable
	{Produce premium notebooks; Produce canvas bags}
	{Strong demand and stable material prices; Weak demand and higher material prices}
	{Market conditions revealed after the product-line decision}
	{Profit in thousand USD for each product-line--market-state pair}

	The only revenue component is an alternative-based variable component:
	\[
	R(A_i,S_j)=q(A_i,S_j)r(A_i).
	\]
	The only cost component is a state-based variable component:
	\[
	C(A_i,S_j)=q(A_i,S_j)c(S_j).
	\]
	Therefore,
	\[
	P(A_i,S_j)
	=
	q(A_i,S_j)r(A_i)
	-
	q(A_i,S_j)c(S_j).
	\]

	\secInicio[sec:problem5-custom-printing-product-line-c1-c3][sec:one-component-models]

	\subsubsection*{C3 - Calculating profit and building the payoff table.}
	\phantomsection
	\label{subsec:problem5-custom-printing-product-line-c3}

	For every alternative--state pair, calculate revenue and cost independently:
	\[
	\begin{aligned}
		R(A_i,S_j)&=q(A_i,S_j)r(A_i),
		\\
		C(A_i,S_j)&=q(A_i,S_j)c(S_j),
		\\
		P(A_i,S_j)&=R(A_i,S_j)-C(A_i,S_j).
	\end{aligned}
	\]

	\begin{center}
		\small
		\begin{tabular}{llll}
			\toprule
			Cell
			& \(R(A_i,S_j)\)
			& \(C(A_i,S_j)\)
			& \(P(A_i,S_j)=R(A_i,S_j)-C(A_i,S_j)\) \\
			\midrule
			\((A_1,S_1)\)
			& \(6(12)=72\)
			& \(6(5)=30\)
			& \(72-30=42\) \\
			
			\((A_1,S_2)\)
			& \(4(12)=48\)
			& \(4(8)=32\)
			& \(48-32=16\) \\
			
			\((A_2,S_1)\)
			& \(5(15)=75\)
			& \(5(5)=25\)
			& \(75-25=50\) \\
			
			\((A_2,S_2)\)
			& \(3(15)=45\)
			& \(3(8)=24\)
			& \(45-24=21\) \\
			\bottomrule
		\end{tabular}
	\end{center}

	The resulting payoff table is:

	\begin{center}
		\begin{tabular}{lrr}
			\toprule
			Alternative
			& Strong demand and stable prices
			& Weak demand and higher prices \\
			\midrule
			Premium notebooks & 42 & 16 \\
			Canvas bags & 50 & 21 \\
			\bottomrule
		\end{tabular}
	\end{center}

	\secInicio[sec:problem5-custom-printing-product-line-c1-c3][sec:one-component-models]

	\subsection{Event Equipment Fleet (C1--C3)}
	\label{sec:problem6-event-equipment-fleet-c1-c3}

	An event-equipment provider must prepare either a compact fleet or an expanded fleet before rental demand is known. Booking enquiries indicate that the compact fleet will fulfil \(6\) rental contracts under high demand and \(3\) under low demand; the expanded fleet has capacity for \(8\) and \(4\) contracts in those respective states.

	Market prices make every completed contract earn \(9\) thousand USD when demand is high and \(6\) thousand USD when it is low, regardless of the fleet selected. Preparing, transporting, and inspecting the equipment requires a fixed outlay of \(22\) thousand USD for the compact fleet or \(34\) thousand USD for the expanded fleet. No other revenue or cost component is modeled, so the provider must compare state-based contract revenue with the chosen fleet's fixed preparation cost.

	\textbf{Questions/Tasks.}
	\begin{enumerate}
		\CriterionTask
		{subsec:problem6-event-equipment-fleet-c1}{1}
		{Interpreting decision alternatives, events, consequences, and states.}

		\CriterionTask
		{subsec:problem6-event-equipment-fleet-c3}{3}
		{Calculating profit and building the payoff table.}
	\end{enumerate}

	\subsubsection*{C1 - Interpreting decision alternatives, events, consequences, and states.}
	\phantomsection
	\label{subsec:problem6-event-equipment-fleet-c1}

	\DecisionElementsTable
	{Prepare a compact fleet; Prepare an expanded fleet}
	{High rental demand; Low rental demand}
	{Rental demand revealed after the fleet-size decision}
	{Profit in thousand USD for each fleet--demand pair}

	The only revenue component is a state-based variable component:
	\[
	R(A_i,S_j)=q(A_i,S_j)r(S_j).
	\]
	The only cost component is an alternative-based fixed component:
	\[
	C(A_i,S_j)=f_c(A_i).
	\]
	Therefore,
	\[
	P(A_i,S_j)
	=
	q(A_i,S_j)r(S_j)
	-
	f_c(A_i).
	\]

	\secInicio[sec:problem6-event-equipment-fleet-c1-c3][sec:one-component-models]

	\subsubsection*{C3 - Calculating profit and building the payoff table.}
	\phantomsection
	\label{subsec:problem6-event-equipment-fleet-c3}

	For every alternative--state pair, calculate
	\[
	\begin{aligned}
		R(A_i,S_j)&=q(A_i,S_j)r(S_j),
		\\
		C(A_i,S_j)&=f_c(A_i),
		\\
		P(A_i,S_j)&=R(A_i,S_j)-C(A_i,S_j).
	\end{aligned}
	\]

	\begin{center}
		\small
		\begin{tabular}{llll}
			\toprule
			Cell
			& \(R(A_i,S_j)\)
			& \(C(A_i,S_j)\)
			& \(P(A_i,S_j)=R(A_i,S_j)-C(A_i,S_j)\) \\
			\midrule
			\((A_1,S_1)\)
			& \(6(9)=54\)
			& \(22\)
			& \(54-22=32\) \\
			
			\((A_1,S_2)\)
			& \(3(6)=18\)
			& \(22\)
			& \(18-22=-4\) \\
			
			\((A_2,S_1)\)
			& \(8(9)=72\)
			& \(34\)
			& \(72-34=38\) \\
			
			\((A_2,S_2)\)
			& \(4(6)=24\)
			& \(34\)
			& \(24-34=-10\) \\
			\bottomrule
		\end{tabular}
	\end{center}

	The resulting payoff table is:

	\begin{center}
		\begin{tabular}{lrr}
			\toprule
			Alternative & High rental demand & Low rental demand \\
			\midrule
			Compact fleet & 32 & \(-4\) \\
			Expanded fleet & 38 & \(-10\) \\
			\bottomrule
		\end{tabular}
	\end{center}

	The negative payoffs in the low-demand state represent losses because the
	fixed preparation cost is greater than the revenue earned.

	\secInicio[sec:problem6-event-equipment-fleet-c1-c3][sec:one-component-models]

	\subsection{Community Arts Sponsorship (C1--C3)}
	\label{sec:problem7-community-arts-sponsorship-c1-c3}

	A community arts organisation must accept either a standard or premium sponsorship agreement before officials confirm whether its event faces routine permit requirements or enhanced safety requirements. The sponsor guarantees a fixed payment of \(52\) thousand USD under the standard agreement or \(68\) thousand USD under the premium agreement, so this revenue depends only on the chosen agreement.

	The authorities' requirements determine the event's fixed regulatory and safety bill: routine permits cost \(18\) thousand USD, while enhanced safety measures cost \(31\) thousand USD, whichever sponsorship agreement was selected. Because neither payment changes with attendance or output, no quantity is required; the organisation simply subtracts the relevant state-based fixed cost from its agreement-based fixed revenue.

	\textbf{Questions/Tasks.}
	\begin{enumerate}
		\CriterionTask
		{subsec:problem7-community-arts-sponsorship-c1}{1}
		{Interpreting decision alternatives, events, consequences, and states.}

		\CriterionTask
		{subsec:problem7-community-arts-sponsorship-c3}{3}
		{Calculating profit and building the payoff table.}
	\end{enumerate}

	\subsubsection*{C1 - Interpreting decision alternatives, events, consequences, and states.}
	\phantomsection
	\label{subsec:problem7-community-arts-sponsorship-c1}

	\DecisionElementsTable
	{Accept the standard sponsorship agreement; Accept the premium sponsorship agreement}
	{Routine permit requirements; Enhanced safety requirements}
	{Event requirements confirmed after the sponsorship decision}
	{Profit in thousand USD for each sponsorship--requirement pair}

	The only revenue component is an alternative-based fixed component:
	\[
	R(A_i,S_j)=f_r(A_i).
	\]
	The only cost component is a state-based fixed component:
	\[
	C(A_i,S_j)=f_c(S_j).
	\]
	Therefore,
	\[
	P(A_i,S_j)
	=
	f_r(A_i)
	-
	f_c(S_j).
	\]

	\secInicio[sec:problem7-community-arts-sponsorship-c1-c3][sec:one-component-models]

	\subsubsection*{C3 - Calculating profit and building the payoff table.}
	\phantomsection
	\label{subsec:problem7-community-arts-sponsorship-c3}

	For every alternative--state pair, calculate
	\[
	\begin{aligned}
		R(A_i,S_j)&=f_r(A_i),
		\\
		C(A_i,S_j)&=f_c(S_j),
		\\
		P(A_i,S_j)&=R(A_i,S_j)-C(A_i,S_j).
	\end{aligned}
	\]

	\begin{center}
		\small
		\begin{tabular}{lrrr}
			\toprule
			Cell
			& \(R(A_i,S_j)\)
			& \(C(A_i,S_j)\)
			& \(P(A_i,S_j)=R(A_i,S_j)-C(A_i,S_j)\) \\
			\midrule
			\((A_1,S_1)\) & 52 & 18 & \(52-18=34\) \\
			\((A_1,S_2)\) & 52 & 31 & \(52-31=21\) \\
			\((A_2,S_1)\) & 68 & 18 & \(68-18=50\) \\
			\((A_2,S_2)\) & 68 & 31 & \(68-31=37\) \\
			\bottomrule
		\end{tabular}
	\end{center}

	The resulting payoff table is:

	\begin{center}
		\begin{tabular}{lrr}
			\toprule
			Alternative
			& Routine requirements
			& Enhanced requirements \\
			\midrule
			Standard sponsorship & 34 & 21 \\
			Premium sponsorship & 50 & 37 \\
			\bottomrule
		\end{tabular}
	\end{center}

	\secInicio[sec:problem7-community-arts-sponsorship-c1-c3][sec:one-component-models]

	\subsection{Distribution Contract Fulfilment (C1--C3)}
	\label{sec:problem8-distribution-contract-fulfilment-c1-c3}

	A distribution company must choose manual or automated packing before a new contract's volume is confirmed as standard or high. Capacity forecasts show manual packing processing \(4\) hundreds of orders under the standard-volume contract and \(7\) under the high-volume contract, while automated packing processes \(5\) and \(8\) hundreds, respectively.

	The customer pays a fixed amount for the realised contract size: \(70\) thousand USD for standard volume or \(104\) thousand USD for high volume, regardless of the fulfilment method. Labour makes manual processing cost \(7\) thousand USD per hundred orders, whereas automation lowers that method-based rate to \(5\) thousand USD per hundred orders. With no other components, the company must deduct the chosen method's quantity-based processing cost from the state's fixed contract payment.

	\textbf{Questions/Tasks.}
	\begin{enumerate}
		\CriterionTask
		{subsec:problem8-distribution-contract-fulfilment-c1}{1}
		{Interpreting decision alternatives, events, consequences, and states.}

		\CriterionTask
		{subsec:problem8-distribution-contract-fulfilment-c3}{3}
		{Calculating profit and building the payoff table.}
	\end{enumerate}

	\subsubsection*{C1 - Interpreting decision alternatives, events, consequences, and states.}
	\phantomsection
	\label{subsec:problem8-distribution-contract-fulfilment-c1}

	\DecisionElementsTable
	{Use manual packing; Use automated packing}
	{Standard-volume contract; High-volume contract}
	{Contract volume revealed after the fulfilment-method decision}
	{Profit in thousand USD for each fulfilment-method--volume pair}

	The only revenue component is a state-based fixed component:
	\[
	R(A_i,S_j)=f_r(S_j).
	\]
	The only cost component is an alternative-based variable component:
	\[
	C(A_i,S_j)=q(A_i,S_j)c(A_i).
	\]
	Therefore,
	\[
	P(A_i,S_j)
	=
	f_r(S_j)
	-
	q(A_i,S_j)c(A_i).
	\]

	\secInicio[sec:problem8-distribution-contract-fulfilment-c1-c3][sec:one-component-models]

	\subsubsection*{C3 - Calculating profit and building the payoff table.}
	\phantomsection
	\label{subsec:problem8-distribution-contract-fulfilment-c3}

	For every alternative--state pair, calculate
	\[
	\begin{aligned}
		R(A_i,S_j)&=f_r(S_j),
		\\
		C(A_i,S_j)&=q(A_i,S_j)c(A_i),
		\\
		P(A_i,S_j)&=R(A_i,S_j)-C(A_i,S_j).
	\end{aligned}
	\]

	\begin{center}
		\small
		\begin{tabular}{llll}
			\toprule
			Cell
			& \(R(A_i,S_j)\)
			& \(C(A_i,S_j)\)
			& \(P(A_i,S_j)=R(A_i,S_j)-C(A_i,S_j)\) \\
			\midrule
			\((A_1,S_1)\)
			& \(70\)
			& \(4(7)=28\)
			& \(70-28=42\) \\
			
			\((A_1,S_2)\)
			& \(104\)
			& \(7(7)=49\)
			& \(104-49=55\) \\
			
			\((A_2,S_1)\)
			& \(70\)
			& \(5(5)=25\)
			& \(70-25=45\) \\
			
			\((A_2,S_2)\)
			& \(104\)
			& \(8(5)=40\)
			& \(104-40=64\) \\
			\bottomrule
		\end{tabular}
	\end{center}

	The resulting payoff table is:

	\begin{center}
		\begin{tabular}{lrr}
			\toprule
			Alternative
			& Standard-volume contract
			& High-volume contract \\
			\midrule
			Manual packing & 42 & 55 \\
			Automated packing & 45 & 64 \\
			\bottomrule
		\end{tabular}
	\end{center}

	\secInicio[sec:problem8-distribution-contract-fulfilment-c1-c3][sec:one-component-models]

	\section{Two-Component Models}
	\label{sec:two-component-models}

	\SectionProblemsTableOfContents{
		\SectionProblemLink{sec:problem9-farm-produce-subscription-c1-c3}{Farm Produce Subscription Plans (C1--C3)} \\
		\SectionProblemLink{sec:problem10-mobile-catering-configuration-c1-c3}{Mobile Catering Configuration (C1--C3)} \\
		\SectionProblemLink{sec:problem11-online-training-programme-c1-c3}{Online Training Programme (C1--C3)} \\
		\SectionProblemLink{sec:problem12-community-venue-membership-c1-c3}{Community Venue Membership (C1--C3)} \\
		\SectionProblemLink{sec:problem13-urban-delivery-fleet-c1-c3}{Urban Delivery Fleet (C1--C3)} \\
		\SectionProblemLink{sec:problem14-data-service-contract-c1-c3}{Data-Service Contract (C1--C3)}
	}

	\subsection{Farm Produce Subscription Plans (C1--C3)}
	\label{sec:problem9-farm-produce-subscription-c1-c3}

	A farm cooperative must offer either a standard or premium produce box before seasonal demand is known to be strong or weak. Its subscription forecast gives the standard box \(7\) hundreds of subscriptions in the strong state and \(4\) in the weak state, and the premium box \(6\) and \(3\) hundreds, respectively.

	The standard box's contents generate a plan-based revenue rate of \(8\) thousand USD per hundred subscriptions, while the premium contents generate \(11\). Strong market conditions add \(3\) thousand USD per hundred subscriptions to either plan's revenue; weak conditions add \(1\). Packing inputs cost \(4\) thousand USD per hundred standard subscriptions or \(6\) per hundred premium subscriptions. Finally, seasonal distribution arrangements create a fixed cost of \(9\) thousand USD under strong demand and \(13\) under weak demand. The cooperative must combine the two revenue rates, then deduct both the plan-based variable cost and the state-based fixed distribution cost.

	\textbf{Questions/Tasks.}
	\begin{enumerate}
		\CriterionTask
		{subsec:problem9-farm-produce-subscription-c1}{1}
		{Interpreting decision alternatives, events, consequences, and states.}

		\CriterionTask
		{subsec:problem9-farm-produce-subscription-c3}{3}
		{Calculating profit and building the payoff table.}
	\end{enumerate}

	\subsubsection*{C1 - Interpreting decision alternatives, events, consequences, and states.}
	\phantomsection
	\label{subsec:problem9-farm-produce-subscription-c1}

	\DecisionElementsTable
	{Offer a standard produce box; Offer a premium produce box}
	{Strong seasonal demand; Weak seasonal demand}
	{Seasonal market conditions revealed after the subscription-plan decision}
	{Profit in thousand USD for each plan--market-state pair}

	Revenue contains alternative-based and state-based variable components:
	\[
	R(A_i,S_j)
	=
	q(A_i,S_j)
	\left[
	r(A_i)+r(S_j)
	\right].
	\]

	Cost contains an alternative-based variable component and a state-based
	fixed component:
	\[
	C(A_i,S_j)
	=
	q(A_i,S_j)c(A_i)+f_c(S_j).
	\]

	Therefore,
	\[
	P(A_i,S_j)
	=
	q(A_i,S_j)
	\left[
	r(A_i)+r(S_j)
	\right]
	-
	\left[
	q(A_i,S_j)c(A_i)+f_c(S_j)
	\right].
	\]

	\secInicio[sec:problem9-farm-produce-subscription-c1-c3][sec:two-component-models]

	\subsubsection*{C3 - Calculating profit and building the payoff table.}
	\phantomsection
	\label{subsec:problem9-farm-produce-subscription-c3}

	For every alternative--state pair, calculate
	\[
	\begin{aligned}
		R(A_i,S_j)
		&=
		q(A_i,S_j)
		\left[
		r(A_i)+r(S_j)
		\right],
		\\
		C(A_i,S_j)
		&=
		q(A_i,S_j)c(A_i)+f_c(S_j),
		\\
		P(A_i,S_j)
		&=
		R(A_i,S_j)-C(A_i,S_j).
	\end{aligned}
	\]

	\begin{center}
		\small
		\begin{tabular}{llll}
			\toprule
			Cell
			& \(R(A_i,S_j)\)
			& \(C(A_i,S_j)\)
			& \(P(A_i,S_j)=R(A_i,S_j)-C(A_i,S_j)\) \\
			\midrule
			\((A_1,S_1)\)
			& \(7(8+3)=77\)
			& \(7(4)+9=37\)
			& \(77-37=40\) \\
			
			\((A_1,S_2)\)
			& \(4(8+1)=36\)
			& \(4(4)+13=29\)
			& \(36-29=7\) \\
			
			\((A_2,S_1)\)
			& \(6(11+3)=84\)
			& \(6(6)+9=45\)
			& \(84-45=39\) \\
			
			\((A_2,S_2)\)
			& \(3(11+1)=36\)
			& \(3(6)+13=31\)
			& \(36-31=5\) \\
			\bottomrule
		\end{tabular}
	\end{center}

	The resulting payoff table is:

	\begin{center}
		\begin{tabular}{lrr}
			\toprule
			Alternative & Strong seasonal demand & Weak seasonal demand \\
			\midrule
			Standard produce box & 40 & 7 \\
			Premium produce box & 39 & 5 \\
			\bottomrule
		\end{tabular}
	\end{center}

	\secInicio[sec:problem9-farm-produce-subscription-c1-c3][sec:two-component-models]

	\subsection{Mobile Catering Configuration (C1--C3)}
	\label{sec:problem10-mobile-catering-configuration-c1-c3}

	A mobile-catering company must prepare a compact catering station or a full mobile kitchen before learning whether the event will be busy with stable ingredient prices or quiet with higher ingredient prices. Booking forecasts indicate \(8\) hundreds of service orders for the compact station in the busy state and \(4\) in the quiet state, compared with \(10\) and \(5\) hundreds for the full kitchen.

	Menu sales earn a configuration-based \(7\) thousand USD per hundred orders from the compact station or \(9\) from the full kitchen. Booking agreements add fixed revenue of \(18\) thousand USD for the compact configuration and \(25\) for the full configuration. Ingredients cost \(3\) thousand USD per hundred orders when prices are stable and \(5\) when prices are higher. In addition, setting up the compact station costs a fixed \(20\) thousand USD, while preparing the full kitchen costs \(32\). The company must retain each component's alternative- or state-based role when calculating profit.

	\textbf{Questions/Tasks.}
	\begin{enumerate}
		\CriterionTask
		{subsec:problem10-mobile-catering-configuration-c1}{1}
		{Interpreting decision alternatives, events, consequences, and states.}

		\CriterionTask
		{subsec:problem10-mobile-catering-configuration-c3}{3}
		{Calculating profit and building the payoff table.}
	\end{enumerate}

	\subsubsection*{C1 - Interpreting decision alternatives, events, consequences, and states.}
	\phantomsection
	\label{subsec:problem10-mobile-catering-configuration-c1}

	\DecisionElementsTable
	{Prepare a compact catering station; Prepare a full mobile kitchen}
	{Busy event with stable ingredient prices; Quiet event with higher ingredient prices}
	{Event demand and ingredient-market conditions revealed after the configuration decision}
	{Profit in thousand USD for each configuration--market-state pair}

	Revenue contains alternative-based variable and fixed components:
	\[
	R(A_i,S_j)
	=
	q(A_i,S_j)r(A_i)+f_r(A_i).
	\]

	Cost contains a state-based variable component and an alternative-based
	fixed component:
	\[
	C(A_i,S_j)
	=
	q(A_i,S_j)c(S_j)+f_c(A_i).
	\]

	Therefore,
	\[
	P(A_i,S_j)
	=
	\left[
	q(A_i,S_j)r(A_i)+f_r(A_i)
	\right]
	-
	\left[
	q(A_i,S_j)c(S_j)+f_c(A_i)
	\right].
	\]

	\secInicio[sec:problem10-mobile-catering-configuration-c1-c3][sec:two-component-models]

	\subsubsection*{C3 - Calculating profit and building the payoff table.}
	\phantomsection
	\label{subsec:problem10-mobile-catering-configuration-c3}

	For every alternative--state pair, calculate
	\[
	\begin{aligned}
		R(A_i,S_j)
		&=
		q(A_i,S_j)r(A_i)+f_r(A_i),
		\\
		C(A_i,S_j)
		&=
		q(A_i,S_j)c(S_j)+f_c(A_i),
		\\
		P(A_i,S_j)
		&=
		R(A_i,S_j)-C(A_i,S_j).
	\end{aligned}
	\]

	\begin{center}
		\small
		\begin{tabular}{llll}
			\toprule
			Cell
			& \(R(A_i,S_j)\)
			& \(C(A_i,S_j)\)
			& \(P(A_i,S_j)=R(A_i,S_j)-C(A_i,S_j)\) \\
			\midrule
			\((A_1,S_1)\)
			& \(8(7)+18=74\)
			& \(8(3)+20=44\)
			& \(74-44=30\) \\
			
			\((A_1,S_2)\)
			& \(4(7)+18=46\)
			& \(4(5)+20=40\)
			& \(46-40=6\) \\
			
			\((A_2,S_1)\)
			& \(10(9)+25=115\)
			& \(10(3)+32=62\)
			& \(115-62=53\) \\
			
			\((A_2,S_2)\)
			& \(5(9)+25=70\)
			& \(5(5)+32=57\)
			& \(70-57=13\) \\
			\bottomrule
		\end{tabular}
	\end{center}

	The resulting payoff table is:

	\begin{center}
		\begin{tabular}{lrr}
			\toprule
			Alternative & Busy event & Quiet event \\
			\midrule
			Compact catering station & 30 & 6 \\
			Full mobile kitchen & 53 & 13 \\
			\bottomrule
		\end{tabular}
	\end{center}

	\secInicio[sec:problem10-mobile-catering-configuration-c1-c3][sec:two-component-models]

	\subsection{Online Training Programme (C1--C3)}
	\label{sec:problem11-online-training-programme-c1-c3}

	An online-training provider must offer its foundation or advanced programme before enrolment is known to come from corporate clients or individuals. Past enquiries imply \(5\) enrolment groups for the foundation programme under corporate enrolment and \(3\) under individual enrolment; the advanced programme expects \(4\) and \(2\) groups, respectively.

	Tuition generates \(14\) thousand USD per foundation group or \(20\) thousand USD per advanced group. The enrolment channel also provides a fixed payment of \(12\) thousand USD for corporate enrolment or \(7\) for individual enrolment. Delivering a group costs \(6\) thousand USD in the corporate state and \(8\) in the individual state, while state-specific platform arrangements add fixed costs of \(10\) and \(14\), respectively. The provider must apply both state-based costs to the forecast groups while keeping the fixed amounts independent of quantity.

	\textbf{Questions/Tasks.}
	\begin{enumerate}
		\CriterionTask
		{subsec:problem11-online-training-programme-c1}{1}
		{Interpreting decision alternatives, events, consequences, and states.}

		\CriterionTask
		{subsec:problem11-online-training-programme-c3}{3}
		{Calculating profit and building the payoff table.}
	\end{enumerate}

	\subsubsection*{C1 - Interpreting decision alternatives, events, consequences, and states.}
	\phantomsection
	\label{subsec:problem11-online-training-programme-c1}

	\DecisionElementsTable
	{Offer the foundation programme; Offer the advanced programme}
	{Corporate enrolment; Individual enrolment}
	{Enrolment type revealed after the programme decision}
	{Profit in thousand USD for each programme--enrolment pair}

	Revenue contains an alternative-based variable component and a state-based
	fixed component:
	\[
	R(A_i,S_j)
	=
	q(A_i,S_j)r(A_i)+f_r(S_j).
	\]

	Cost contains state-based variable and fixed components:
	\[
	C(A_i,S_j)
	=
	q(A_i,S_j)c(S_j)+f_c(S_j).
	\]

	Therefore,
	\[
	P(A_i,S_j)
	=
	\left[
	q(A_i,S_j)r(A_i)+f_r(S_j)
	\right]
	-
	\left[
	q(A_i,S_j)c(S_j)+f_c(S_j)
	\right].
	\]

	\secInicio[sec:problem11-online-training-programme-c1-c3][sec:two-component-models]

	\subsubsection*{C3 - Calculating profit and building the payoff table.}
	\phantomsection
	\label{subsec:problem11-online-training-programme-c3}

	For every alternative--state pair, calculate
	\[
	\begin{aligned}
		R(A_i,S_j)
		&=
		q(A_i,S_j)r(A_i)+f_r(S_j),
		\\
		C(A_i,S_j)
		&=
		q(A_i,S_j)c(S_j)+f_c(S_j),
		\\
		P(A_i,S_j)
		&=
		R(A_i,S_j)-C(A_i,S_j).
	\end{aligned}
	\]

	\begin{center}
		\small
		\begin{tabular}{llll}
			\toprule
			Cell
			& \(R(A_i,S_j)\)
			& \(C(A_i,S_j)\)
			& \(P(A_i,S_j)=R(A_i,S_j)-C(A_i,S_j)\) \\
			\midrule
			\((A_1,S_1)\)
			& \(5(14)+12=82\)
			& \(5(6)+10=40\)
			& \(82-40=42\) \\
			
			\((A_1,S_2)\)
			& \(3(14)+7=49\)
			& \(3(8)+14=38\)
			& \(49-38=11\) \\
			
			\((A_2,S_1)\)
			& \(4(20)+12=92\)
			& \(4(6)+10=34\)
			& \(92-34=58\) \\
			
			\((A_2,S_2)\)
			& \(2(20)+7=47\)
			& \(2(8)+14=30\)
			& \(47-30=17\) \\
			\bottomrule
		\end{tabular}
	\end{center}

	The resulting payoff table is:

	\begin{center}
		\begin{tabular}{lrr}
			\toprule
			Alternative & Corporate enrolment & Individual enrolment \\
			\midrule
			Foundation programme & 42 & 11 \\
			Advanced programme & 58 & 17 \\
			\bottomrule
		\end{tabular}
	\end{center}

	\secInicio[sec:problem11-online-training-programme-c1-c3][sec:two-component-models]

	\subsection{Community Venue Membership (C1--C3)}
	\label{sec:problem12-community-venue-membership-c1-c3}

	A community organisation must select standard or expanded venue membership before booking demand is known to be high or low. Reservations are forecast at \(6\) bookings for standard membership under high demand and \(3\) under low demand, while expanded membership enables \(8\) and \(4\) bookings, respectively.

	The market supports revenue of \(10\) thousand USD per booking in the high-demand state and \(7\) in the low-demand state. A sponsor also contributes a fixed \(15\) thousand USD when the standard plan is selected or \(22\) for the expanded plan. Annual membership itself costs a fixed \(28\) thousand USD for the standard plan and \(40\) for the expanded plan, while permits add a state-based fixed cost of \(8\) under high demand or \(13\) under low demand. The organisation must calculate revenue and subtract both fixed cost sources without multiplying them by bookings.

	\textbf{Questions/Tasks.}
	\begin{enumerate}
		\CriterionTask
		{subsec:problem12-community-venue-membership-c1}{1}
		{Interpreting decision alternatives, events, consequences, and states.}

		\CriterionTask
		{subsec:problem12-community-venue-membership-c3}{3}
		{Calculating profit and building the payoff table.}
	\end{enumerate}

	\subsubsection*{C1 - Interpreting decision alternatives, events, consequences, and states.}
	\phantomsection
	\label{subsec:problem12-community-venue-membership-c1}

	\DecisionElementsTable
	{Select the standard venue membership; Select the expanded venue membership}
	{High booking demand; Low booking demand}
	{Booking demand revealed after the membership-plan decision}
	{Profit in thousand USD for each membership--demand pair}

	Revenue contains a state-based variable component and an alternative-based
	fixed component:
	\[
	R(A_i,S_j)
	=
	q(A_i,S_j)r(S_j)+f_r(A_i).
	\]

	Cost contains alternative-based and state-based fixed components:
	\[
	C(A_i,S_j)
	=
	f_c(A_i)+f_c(S_j).
	\]

	Therefore,
	\[
	P(A_i,S_j)
	=
	\left[
	q(A_i,S_j)r(S_j)+f_r(A_i)
	\right]
	-
	\left[
	f_c(A_i)+f_c(S_j)
	\right].
	\]

	\secInicio[sec:problem12-community-venue-membership-c1-c3][sec:two-component-models]

	\subsubsection*{C3 - Calculating profit and building the payoff table.}
	\phantomsection
	\label{subsec:problem12-community-venue-membership-c3}

	For every alternative--state pair, calculate
	\[
	\begin{aligned}
		R(A_i,S_j)
		&=
		q(A_i,S_j)r(S_j)+f_r(A_i),
		\\
		C(A_i,S_j)
		&=
		f_c(A_i)+f_c(S_j),
		\\
		P(A_i,S_j)
		&=
		R(A_i,S_j)-C(A_i,S_j).
	\end{aligned}
	\]

	\begin{center}
		\small
		\begin{tabular}{llll}
			\toprule
			Cell
			& \(R(A_i,S_j)\)
			& \(C(A_i,S_j)\)
			& \(P(A_i,S_j)=R(A_i,S_j)-C(A_i,S_j)\) \\
			\midrule
			\((A_1,S_1)\)
			& \(6(10)+15=75\)
			& \(28+8=36\)
			& \(75-36=39\) \\
			
			\((A_1,S_2)\)
			& \(3(7)+15=36\)
			& \(28+13=41\)
			& \(36-41=-5\) \\
			
			\((A_2,S_1)\)
			& \(8(10)+22=102\)
			& \(40+8=48\)
			& \(102-48=54\) \\
			
			\((A_2,S_2)\)
			& \(4(7)+22=50\)
			& \(40+13=53\)
			& \(50-53=-3\) \\
			\bottomrule
		\end{tabular}
	\end{center}

	The resulting payoff table is:

	\begin{center}
		\begin{tabular}{lrr}
			\toprule
			Alternative & High booking demand & Low booking demand \\
			\midrule
			Standard venue membership & 39 & \(-5\) \\
			Expanded venue membership & 54 & \(-3\) \\
			\bottomrule
		\end{tabular}
	\end{center}

	The negative payoffs under low booking demand represent losses because the
	combined fixed costs exceed the revenue earned.

	\secInicio[sec:problem12-community-venue-membership-c1-c3][sec:two-component-models]

	\subsection{Urban Delivery Fleet (C1--C3)}
	\label{sec:problem13-urban-delivery-fleet-c1-c3}

	An urban-delivery company must deploy conventional vans or electric cargo bikes before conditions resolve into high demand with light traffic or low demand with heavy traffic. Route forecasts assign the vans \(9\) delivery batches in the first state and \(5\) in the second, while cargo bikes can complete \(12\) and \(6\) batches, respectively.

	Customer contracts pay \(8\) thousand USD per batch in the high-demand state and \(6\) in the low-demand state. Meeting state targets also brings a fixed performance payment of \(14\) thousand USD with high demand or \(8\) with low demand. Van operation costs \(3\) thousand USD per batch and cargo-bike operation \(2\); traffic then adds \(1\) thousand USD per batch in light traffic or \(3\) in heavy traffic. The company must combine the fleet-based and traffic-based rates for every delivered batch before deducting costs.

	\textbf{Questions/Tasks.}
	\begin{enumerate}
		\CriterionTask
		{subsec:problem13-urban-delivery-fleet-c1}{1}
		{Interpreting decision alternatives, events, consequences, and states.}

		\CriterionTask
		{subsec:problem13-urban-delivery-fleet-c3}{3}
		{Calculating profit and building the payoff table.}
	\end{enumerate}

	\subsubsection*{C1 - Interpreting decision alternatives, events, consequences, and states.}
	\phantomsection
	\label{subsec:problem13-urban-delivery-fleet-c1}

	\DecisionElementsTable
	{Use a conventional van fleet; Use an electric cargo-bike fleet}
	{High demand with light traffic; Low demand with heavy traffic}
	{Delivery demand and traffic conditions revealed after the fleet decision}
	{Profit in thousand USD for each fleet--market-state pair}

	Revenue contains state-based variable and fixed components:
	\[
	R(A_i,S_j)
	=
	q(A_i,S_j)r(S_j)+f_r(S_j).
	\]

	Cost contains alternative-based and state-based variable components:
	\[
	C(A_i,S_j)
	=
	q(A_i,S_j)
	\left[
	c(A_i)+c(S_j)
	\right].
	\]

	Therefore,
	\[
	P(A_i,S_j)
	=
	\left[
	q(A_i,S_j)r(S_j)+f_r(S_j)
	\right]
	-
	q(A_i,S_j)
	\left[
	c(A_i)+c(S_j)
	\right].
	\]

	\secInicio[sec:problem13-urban-delivery-fleet-c1-c3][sec:two-component-models]

	\subsubsection*{C3 - Calculating profit and building the payoff table.}
	\phantomsection
	\label{subsec:problem13-urban-delivery-fleet-c3}

	For every alternative--state pair, calculate
	\[
	\begin{aligned}
		R(A_i,S_j)
		&=
		q(A_i,S_j)r(S_j)+f_r(S_j),
		\\
		C(A_i,S_j)
		&=
		q(A_i,S_j)
		\left[
		c(A_i)+c(S_j)
		\right],
		\\
		P(A_i,S_j)
		&=
		R(A_i,S_j)-C(A_i,S_j).
	\end{aligned}
	\]

	\begin{center}
		\small
		\begin{tabular}{llll}
			\toprule
			Cell
			& \(R(A_i,S_j)\)
			& \(C(A_i,S_j)\)
			& \(P(A_i,S_j)=R(A_i,S_j)-C(A_i,S_j)\) \\
			\midrule
			\((A_1,S_1)\)
			& \(9(8)+14=86\)
			& \(9(3+1)=36\)
			& \(86-36=50\) \\
			
			\((A_1,S_2)\)
			& \(5(6)+8=38\)
			& \(5(3+3)=30\)
			& \(38-30=8\) \\
			
			\((A_2,S_1)\)
			& \(12(8)+14=110\)
			& \(12(2+1)=36\)
			& \(110-36=74\) \\
			
			\((A_2,S_2)\)
			& \(6(6)+8=44\)
			& \(6(2+3)=30\)
			& \(44-30=14\) \\
			\bottomrule
		\end{tabular}
	\end{center}

	The resulting payoff table is:

	\begin{center}
		\begin{tabular}{lrr}
			\toprule
			Alternative & High demand, light traffic & Low demand, heavy traffic \\
			\midrule
			Conventional van fleet & 50 & 8 \\
			Electric cargo-bike fleet & 74 & 14 \\
			\bottomrule
		\end{tabular}
	\end{center}

	\secInicio[sec:problem13-urban-delivery-fleet-c1-c3][sec:two-component-models]

	\subsection{Data-Service Contract (C1--C3)}
	\label{sec:problem14-data-service-contract-c1-c3}

	A data-services company must choose standard or automated processing before a contract is confirmed as standard-volume or high-volume. Work estimates assign standard processing \(4\) batches under standard volume and \(7\) under high volume, while automation handles \(5\) and \(9\) batches, respectively.

	The method agreement pays fixed revenue of \(55\) thousand USD for standard processing or \(72\) for automation. The client adds a service-level payment of \(16\) thousand USD in the standard-volume state and \(9\) in the high-volume state. Processing costs \(6\) thousand USD per batch with the standard method or \(4\) with automation, and implementation adds a method-based fixed cost of \(18\) thousand USD or \(31\), respectively. The company must keep both revenue payments fixed while applying only the processing rate to batch quantity.

	\textbf{Questions/Tasks.}
	\begin{enumerate}
		\CriterionTask
		{subsec:problem14-data-service-contract-c1}{1}
		{Interpreting decision alternatives, events, consequences, and states.}

		\CriterionTask
		{subsec:problem14-data-service-contract-c3}{3}
		{Calculating profit and building the payoff table.}
	\end{enumerate}

	\subsubsection*{C1 - Interpreting decision alternatives, events, consequences, and states.}
	\phantomsection
	\label{subsec:problem14-data-service-contract-c1}

	\DecisionElementsTable
	{Use standard data processing; Use automated data processing}
	{Standard-volume contract; High-volume contract}
	{Contract requirements confirmed after the processing-method decision}
	{Profit in thousand USD for each processing-method--contract-state pair}

	Revenue contains alternative-based and state-based fixed components:
	\[
	R(A_i,S_j)
	=
	f_r(A_i)+f_r(S_j).
	\]

	Cost contains alternative-based variable and fixed components:
	\[
	C(A_i,S_j)
	=
	q(A_i,S_j)c(A_i)+f_c(A_i).
	\]

	Therefore,
	\[
	P(A_i,S_j)
	=
	\left[
	f_r(A_i)+f_r(S_j)
	\right]
	-
	\left[
	q(A_i,S_j)c(A_i)+f_c(A_i)
	\right].
	\]

	\secInicio[sec:problem14-data-service-contract-c1-c3][sec:two-component-models]

	\subsubsection*{C3 - Calculating profit and building the payoff table.}
	\phantomsection
	\label{subsec:problem14-data-service-contract-c3}

	For every alternative--state pair, calculate
	\[
	\begin{aligned}
		R(A_i,S_j)
		&=
		f_r(A_i)+f_r(S_j),
		\\
		C(A_i,S_j)
		&=
		q(A_i,S_j)c(A_i)+f_c(A_i),
		\\
		P(A_i,S_j)
		&=
		R(A_i,S_j)-C(A_i,S_j).
	\end{aligned}
	\]

	\begin{center}
		\small
		\begin{tabular}{llll}
			\toprule
			Cell
			& \(R(A_i,S_j)\)
			& \(C(A_i,S_j)\)
			& \(P(A_i,S_j)=R(A_i,S_j)-C(A_i,S_j)\) \\
			\midrule
			\((A_1,S_1)\)
			& \(55+16=71\)
			& \(4(6)+18=42\)
			& \(71-42=29\) \\
			
			\((A_1,S_2)\)
			& \(55+9=64\)
			& \(7(6)+18=60\)
			& \(64-60=4\) \\
			
			\((A_2,S_1)\)
			& \(72+16=88\)
			& \(5(4)+31=51\)
			& \(88-51=37\) \\
			
			\((A_2,S_2)\)
			& \(72+9=81\)
			& \(9(4)+31=67\)
			& \(81-67=14\) \\
			\bottomrule
		\end{tabular}
	\end{center}

	The resulting payoff table is:

	\begin{center}
		\begin{tabular}{lrr}
			\toprule
			Alternative & Standard-volume contract & High-volume contract \\
			\midrule
			Standard data processing & 29 & 4 \\
			Automated data processing & 37 & 14 \\
			\bottomrule
		\end{tabular}
	\end{center}

	\secInicio[sec:problem14-data-service-contract-c1-c3][sec:two-component-models]

	\section{Three-Component Models}
	\label{sec:three-component-models}

	\SectionProblemsTableOfContents{
		\SectionProblemLink{sec:problem15-regional-parcel-delivery-c1-c3}{Regional Parcel Delivery (C1--C3)} \\
		\SectionProblemLink{sec:problem16-solar-installation-programme-c1-c3}{Solar Installation Programme (C1--C3)} \\
		\SectionProblemLink{sec:problem17-hybrid-conference-production-c1-c3}{Hybrid Conference Production (C1--C3)} \\
		\SectionProblemLink{sec:problem18-cold-storage-export-contract-c1-c3}{Cold-Storage Export Contract (C1--C3)}
	}

	\subsection{Regional Parcel Delivery (C1--C3)}
	\label{sec:problem15-regional-parcel-delivery-c1-c3}

	A regional logistics company must select its standard or express delivery plan before conditions become high demand with light traffic or moderate demand with heavy traffic. Forecasts give the standard plan \(8\) hundreds of delivery batches in the first state and \(5\) in the second, and the express plan \(10\) and \(6\) hundreds, respectively.

	The standard service earns \(9\) thousand USD per hundred batches and express service earns \(11\); demand conditions add \(3\) per hundred batches under high demand or \(2\) under moderate demand. Service agreements also pay a fixed \(16\) thousand USD for the standard plan and \(24\) for express. Operating costs are \(4\) thousand USD per hundred batches for standard and \(5\) for express, while traffic adds \(1\) per hundred batches when light and \(3\) when heavy. Distribution arrangements finally cost a fixed \(10\) thousand USD in the high-demand, light-traffic state and \(15\) in the moderate-demand, heavy-traffic state.

	\textbf{Questions/Tasks.}
	\begin{enumerate}
		\CriterionTask
		{subsec:problem15-regional-parcel-delivery-c1}{1}
		{Interpreting decision alternatives, events, consequences, and states.}

		\CriterionTask
		{subsec:problem15-regional-parcel-delivery-c3}{3}
		{Calculating profit and building the payoff table.}
	\end{enumerate}

	\subsubsection*{C1 - Interpreting decision alternatives, events, consequences, and states.}
	\phantomsection
	\label{subsec:problem15-regional-parcel-delivery-c1}

	\DecisionElementsTable
	{Use the standard delivery plan; Use the express delivery plan}
	{High demand with light traffic; Moderate demand with heavy traffic}
	{Demand and traffic conditions revealed after the delivery-plan decision}
	{Profit in thousand USD for each delivery-plan--state pair}

	Revenue contains alternative-based and state-based variable components and
	an alternative-based fixed component:
	\[
	R(A_i,S_j)
	=
	q(A_i,S_j)
	\left[
	r(A_i)+r(S_j)
	\right]
	+
	f_r(A_i).
	\]

	Cost contains alternative-based and state-based variable components and a
	state-based fixed component:
	\[
	C(A_i,S_j)
	=
	q(A_i,S_j)
	\left[
	c(A_i)+c(S_j)
	\right]
	+
	f_c(S_j).
	\]

	Therefore,
	\[
	\begin{aligned}
		P(A_i,S_j)
		={}&
		\left\{
		q(A_i,S_j)
		\left[
		r(A_i)+r(S_j)
		\right]
		+
		f_r(A_i)
		\right\}
		\\
		&-
		\left\{
		q(A_i,S_j)
		\left[
		c(A_i)+c(S_j)
		\right]
		+
		f_c(S_j)
		\right\}.
	\end{aligned}
	\]

	\secInicio[sec:problem15-regional-parcel-delivery-c1-c3][sec:three-component-models]

	\subsubsection*{C3 - Calculating profit and building the payoff table.}
	\phantomsection
	\label{subsec:problem15-regional-parcel-delivery-c3}

	For every alternative--state pair, calculate
	\[
	\begin{aligned}
		R(A_i,S_j)
		&=
		q(A_i,S_j)
		\left[
		r(A_i)+r(S_j)
		\right]
		+
		f_r(A_i),
		\\
		C(A_i,S_j)
		&=
		q(A_i,S_j)
		\left[
		c(A_i)+c(S_j)
		\right]
		+
		f_c(S_j),
		\\
		P(A_i,S_j)
		&=
		R(A_i,S_j)-C(A_i,S_j).
	\end{aligned}
	\]

	\begin{center}
		\small
		\begin{tabular}{llll}
			\toprule
			Cell
			& \(R(A_i,S_j)\)
			& \(C(A_i,S_j)\)
			& \(P(A_i,S_j)=R(A_i,S_j)-C(A_i,S_j)\) \\
			\midrule
			\((A_1,S_1)\)
			& \(8(9+3)+16=112\)
			& \(8(4+1)+10=50\)
			& \(112-50=62\) \\
			
			\((A_1,S_2)\)
			& \(5(9+2)+16=71\)
			& \(5(4+3)+15=50\)
			& \(71-50=21\) \\
			
			\((A_2,S_1)\)
			& \(10(11+3)+24=164\)
			& \(10(5+1)+10=70\)
			& \(164-70=94\) \\
			
			\((A_2,S_2)\)
			& \(6(11+2)+24=102\)
			& \(6(5+3)+15=63\)
			& \(102-63=39\) \\
			\bottomrule
		\end{tabular}
	\end{center}

	The resulting payoff table is:

	\begin{center}
		\begin{tabular}{lrr}
			\toprule
			Alternative
			& High demand, light traffic
			& Moderate demand, heavy traffic \\
			\midrule
			Standard delivery plan & 62 & 21 \\
			Express delivery plan & 94 & 39 \\
			\bottomrule
		\end{tabular}
	\end{center}

	\secInicio[sec:problem15-regional-parcel-delivery-c1-c3][sec:three-component-models]

	\subsection{Solar Installation Programme (C1--C3)}
	\label{sec:problem16-solar-installation-programme-c1-c3}

	A renewable-energy company must offer a standard-panel or high-efficiency-panel programme before demand and permits resolve into strong commercial demand or limited residential demand. It forecasts \(7\) installation batches for standard panels in the strong state and \(4\) in the limited state, compared with \(6\) and \(3\) batches for high-efficiency panels.

	Standard panels earn \(12\) thousand USD per batch and high-efficiency panels \(18\), while the market adds \(4\) per batch in the strong state or \(2\) in the limited state. An incentive provides fixed revenue of \(20\) thousand USD under strong commercial demand and \(11\) under limited residential demand. Installation inputs cost \(6\) thousand USD per standard-panel batch or \(8\) per high-efficiency batch. Programme setup costs a fixed \(18\) thousand USD for standard panels or \(25\) for high-efficiency panels, and permits add fixed costs of \(9\) in the strong state or \(14\) in the limited state.

	\textbf{Questions/Tasks.}
	\begin{enumerate}
		\CriterionTask
		{subsec:problem16-solar-installation-programme-c1}{1}
		{Interpreting decision alternatives, events, consequences, and states.}

		\CriterionTask
		{subsec:problem16-solar-installation-programme-c3}{3}
		{Calculating profit and building the payoff table.}
	\end{enumerate}

	\subsubsection*{C1 - Interpreting decision alternatives, events, consequences, and states.}
	\phantomsection
	\label{subsec:problem16-solar-installation-programme-c1}

	\DecisionElementsTable
	{Offer the standard-panel programme; Offer the high-efficiency-panel programme}
	{Strong commercial demand; Limited residential demand}
	{Customer demand and permit conditions revealed after the programme decision}
	{Profit in thousand USD for each programme--market-state pair}

	Revenue contains alternative-based and state-based variable components and
	a state-based fixed component:
	\[
	R(A_i,S_j)
	=
	q(A_i,S_j)
	\left[
	r(A_i)+r(S_j)
	\right]
	+
	f_r(S_j).
	\]

	Cost contains an alternative-based variable component and alternative-based
	and state-based fixed components:
	\[
	C(A_i,S_j)
	=
	q(A_i,S_j)c(A_i)
	+
	f_c(A_i)
	+
	f_c(S_j).
	\]

	Therefore,
	\[
	\begin{aligned}
		P(A_i,S_j)
		={}&
		\left\{
		q(A_i,S_j)
		\left[
		r(A_i)+r(S_j)
		\right]
		+
		f_r(S_j)
		\right\}
		\\
		&-
		\left\{
		q(A_i,S_j)c(A_i)
		+
		f_c(A_i)
		+
		f_c(S_j)
		\right\}.
	\end{aligned}
	\]

	\secInicio[sec:problem16-solar-installation-programme-c1-c3][sec:three-component-models]

	\subsubsection*{C3 - Calculating profit and building the payoff table.}
	\phantomsection
	\label{subsec:problem16-solar-installation-programme-c3}

	For every alternative--state pair, calculate
	\[
	\begin{aligned}
		R(A_i,S_j)
		&=
		q(A_i,S_j)
		\left[
		r(A_i)+r(S_j)
		\right]
		+
		f_r(S_j),
		\\
		C(A_i,S_j)
		&=
		q(A_i,S_j)c(A_i)
		+
		f_c(A_i)
		+
		f_c(S_j),
		\\
		P(A_i,S_j)
		&=
		R(A_i,S_j)-C(A_i,S_j).
	\end{aligned}
	\]

	\begin{center}
		\small
		\begin{tabular}{llll}
			\toprule
			Cell
			& \(R(A_i,S_j)\)
			& \(C(A_i,S_j)\)
			& \(P(A_i,S_j)=R(A_i,S_j)-C(A_i,S_j)\) \\
			\midrule
			\((A_1,S_1)\)
			& \(7(12+4)+20=132\)
			& \(7(6)+18+9=69\)
			& \(132-69=63\) \\
			
			\((A_1,S_2)\)
			& \(4(12+2)+11=67\)
			& \(4(6)+18+14=56\)
			& \(67-56=11\) \\
			
			\((A_2,S_1)\)
			& \(6(18+4)+20=152\)
			& \(6(8)+25+9=82\)
			& \(152-82=70\) \\
			
			\((A_2,S_2)\)
			& \(3(18+2)+11=71\)
			& \(3(8)+25+14=63\)
			& \(71-63=8\) \\
			\bottomrule
		\end{tabular}
	\end{center}

	The resulting payoff table is:

	\begin{center}
		\begin{tabular}{lrr}
			\toprule
			Alternative
			& Strong commercial demand
			& Limited residential demand \\
			\midrule
			Standard-panel programme & 63 & 11 \\
			High-efficiency-panel programme & 70 & 8 \\
			\bottomrule
		\end{tabular}
	\end{center}

	\secInicio[sec:problem16-solar-installation-programme-c1-c3][sec:three-component-models]

	\subsection{Hybrid Conference Production (C1--C3)}
	\label{sec:problem17-hybrid-conference-production-c1-c3}

	An events company must choose a standard or immersive hybrid-conference package before attendance and sponsorship are known to be strong or limited. Production forecasts allocate \(6\) event segments to the standard package in the strong state and \(4\) in the limited state, while the immersive package uses \(8\) and \(5\) segments, respectively.

	Each standard segment earns \(13\) thousand USD and each immersive segment \(16\). Booking contracts add fixed revenue of \(14\) thousand USD for the standard package or \(22\) for the immersive package, while sponsors add a further fixed \(18\) in the strong state or \(9\) in the limited state. Delivery costs \(5\) thousand USD per segment under strong conditions and \(7\) under limited conditions. Package setup adds a fixed cost of \(24\) for standard or \(37\) for immersive, and licensing adds a state-based fixed cost of \(8\) or \(13\), respectively.

	\textbf{Questions/Tasks.}
	\begin{enumerate}
		\CriterionTask
		{subsec:problem17-hybrid-conference-production-c1}{1}
		{Interpreting decision alternatives, events, consequences, and states.}

		\CriterionTask
		{subsec:problem17-hybrid-conference-production-c3}{3}
		{Calculating profit and building the payoff table.}
	\end{enumerate}

	\subsubsection*{C1 - Interpreting decision alternatives, events, consequences, and states.}
	\phantomsection
	\label{subsec:problem17-hybrid-conference-production-c1}

	\DecisionElementsTable
	{Use the standard hybrid package; Use the immersive hybrid package}
	{Strong attendance and sponsorship; Limited attendance and sponsorship}
	{Attendance and sponsorship conditions revealed after the package decision}
	{Profit in thousand USD for each production-package--state pair}

	Revenue contains an alternative-based variable component and
	alternative-based and state-based fixed components:
	\[
	R(A_i,S_j)
	=
	q(A_i,S_j)r(A_i)
	+
	f_r(A_i)
	+
	f_r(S_j).
	\]

	Cost contains a state-based variable component and alternative-based and
	state-based fixed components:
	\[
	C(A_i,S_j)
	=
	q(A_i,S_j)c(S_j)
	+
	f_c(A_i)
	+
	f_c(S_j).
	\]

	Therefore,
	\[
	\begin{aligned}
		P(A_i,S_j)
		={}&
		\left\{
		q(A_i,S_j)r(A_i)
		+
		f_r(A_i)
		+
		f_r(S_j)
		\right\}
		\\
		&-
		\left\{
		q(A_i,S_j)c(S_j)
		+
		f_c(A_i)
		+
		f_c(S_j)
		\right\}.
	\end{aligned}
	\]

	\secInicio[sec:problem17-hybrid-conference-production-c1-c3][sec:three-component-models]

	\subsubsection*{C3 - Calculating profit and building the payoff table.}
	\phantomsection
	\label{subsec:problem17-hybrid-conference-production-c3}

	For every alternative--state pair, calculate
	\[
	\begin{aligned}
		R(A_i,S_j)
		&=
		q(A_i,S_j)r(A_i)
		+
		f_r(A_i)
		+
		f_r(S_j),
		\\
		C(A_i,S_j)
		&=
		q(A_i,S_j)c(S_j)
		+
		f_c(A_i)
		+
		f_c(S_j),
		\\
		P(A_i,S_j)
		&=
		R(A_i,S_j)-C(A_i,S_j).
	\end{aligned}
	\]

	\begin{center}
		\small
		\begin{tabular}{llll}
			\toprule
			Cell
			& \(R(A_i,S_j)\)
			& \(C(A_i,S_j)\)
			& \(P(A_i,S_j)=R(A_i,S_j)-C(A_i,S_j)\) \\
			\midrule
			\((A_1,S_1)\)
			& \(6(13)+14+18=110\)
			& \(6(5)+24+8=62\)
			& \(110-62=48\) \\
			
			\((A_1,S_2)\)
			& \(4(13)+14+9=75\)
			& \(4(7)+24+13=65\)
			& \(75-65=10\) \\
			
			\((A_2,S_1)\)
			& \(8(16)+22+18=168\)
			& \(8(5)+37+8=85\)
			& \(168-85=83\) \\
			
			\((A_2,S_2)\)
			& \(5(16)+22+9=111\)
			& \(5(7)+37+13=85\)
			& \(111-85=26\) \\
			\bottomrule
		\end{tabular}
	\end{center}

	The resulting payoff table is:

	\begin{center}
		\begin{tabular}{lrr}
			\toprule
			Alternative
			& Strong attendance and sponsorship
			& Limited attendance and sponsorship \\
			\midrule
			Standard hybrid package & 48 & 10 \\
			Immersive hybrid package & 83 & 26 \\
			\bottomrule
		\end{tabular}
	\end{center}

	\secInicio[sec:problem17-hybrid-conference-production-c1-c3][sec:three-component-models]

	\subsection{Cold-Storage Export Contract (C1--C3)}
	\label{sec:problem18-cold-storage-export-contract-c1-c3}

	A food-export company must choose a modular or large cold-storage facility before export demand is strong or limited. Capacity and demand forecasts give the modular facility \(9\) storage batches under strong demand and \(5\) under limited demand, and the large facility \(12\) and \(7\) batches, respectively.

	Export conditions set revenue at \(10\) thousand USD per batch under strong demand or \(7\) under limited demand. The chosen capacity earns fixed revenue of \(20\) thousand USD for the modular facility or \(30\) for the large facility, and the export programme adds fixed revenue of \(16\) in the strong state or \(8\) in the limited state. Operating rates are \(4\) thousand USD per modular batch and \(3\) per large-facility batch; energy adds \(2\) per batch under strong demand or \(4\) under limited demand. Preparing the configuration also costs a fixed \(28\) thousand USD for modular storage or \(42\) for large storage.

	\textbf{Questions/Tasks.}
	\begin{enumerate}
		\CriterionTask
		{subsec:problem18-cold-storage-export-contract-c1}{1}
		{Interpreting decision alternatives, events, consequences, and states.}

		\CriterionTask
		{subsec:problem18-cold-storage-export-contract-c3}{3}
		{Calculating profit and building the payoff table.}
	\end{enumerate}

	\subsubsection*{C1 - Interpreting decision alternatives, events, consequences, and states.}
	\phantomsection
	\label{subsec:problem18-cold-storage-export-contract-c1}

	\DecisionElementsTable
	{Use the modular cold-storage facility; Use the large cold-storage facility}
	{Strong export demand; Limited export demand}
	{Export demand and energy-market conditions revealed after the facility decision}
	{Profit in thousand USD for each storage-facility--state pair}

	Revenue contains a state-based variable component and alternative-based and
	state-based fixed components:
	\[
	R(A_i,S_j)
	=
	q(A_i,S_j)r(S_j)
	+
	f_r(A_i)
	+
	f_r(S_j).
	\]

	Cost contains alternative-based and state-based variable components and an
	alternative-based fixed component:
	\[
	C(A_i,S_j)
	=
	q(A_i,S_j)
	\left[
	c(A_i)+c(S_j)
	\right]
	+
	f_c(A_i).
	\]

	Therefore,
	\[
	\begin{aligned}
		P(A_i,S_j)
		={}&
		\left\{
		q(A_i,S_j)r(S_j)
		+
		f_r(A_i)
		+
		f_r(S_j)
		\right\}
		\\
		&-
		\left\{
		q(A_i,S_j)
		\left[
		c(A_i)+c(S_j)
		\right]
		+
		f_c(A_i)
		\right\}.
	\end{aligned}
	\]

	\secInicio[sec:problem18-cold-storage-export-contract-c1-c3][sec:three-component-models]

	\subsubsection*{C3 - Calculating profit and building the payoff table.}
	\phantomsection
	\label{subsec:problem18-cold-storage-export-contract-c3}

	For every alternative--state pair, calculate
	\[
	\begin{aligned}
		R(A_i,S_j)
		&=
		q(A_i,S_j)r(S_j)
		+
		f_r(A_i)
		+
		f_r(S_j),
		\\
		C(A_i,S_j)
		&=
		q(A_i,S_j)
		\left[
		c(A_i)+c(S_j)
		\right]
		+
		f_c(A_i),
		\\
		P(A_i,S_j)
		&=
		R(A_i,S_j)-C(A_i,S_j).
	\end{aligned}
	\]

	\begin{center}
		\small
		\begin{tabular}{llll}
			\toprule
			Cell
			& \(R(A_i,S_j)\)
			& \(C(A_i,S_j)\)
			& \(P(A_i,S_j)=R(A_i,S_j)-C(A_i,S_j)\) \\
			\midrule
			\((A_1,S_1)\)
			& \(9(10)+20+16=126\)
			& \(9(4+2)+28=82\)
			& \(126-82=44\) \\
			
			\((A_1,S_2)\)
			& \(5(7)+20+8=63\)
			& \(5(4+4)+28=68\)
			& \(63-68=-5\) \\
			
			\((A_2,S_1)\)
			& \(12(10)+30+16=166\)
			& \(12(3+2)+42=102\)
			& \(166-102=64\) \\
			
			\((A_2,S_2)\)
			& \(7(7)+30+8=87\)
			& \(7(3+4)+42=91\)
			& \(87-91=-4\) \\
			\bottomrule
		\end{tabular}
	\end{center}

	The resulting payoff table is:

	\begin{center}
		\begin{tabular}{lrr}
			\toprule
			Alternative & Strong export demand & Limited export demand \\
			\midrule
			Modular cold-storage facility & 44 & \(-5\) \\
			Large cold-storage facility & 64 & \(-4\) \\
			\bottomrule
		\end{tabular}
	\end{center}

	The negative payoffs under limited export demand represent losses because
	the revenue earned does not cover the variable and fixed costs.

	\secInicio[sec:problem18-cold-storage-export-contract-c1-c3][sec:three-component-models]

	\section{Complete Four-Component Models}
	\label{sec:complete-four-component-models}

	\SectionProblemsTableOfContents{
		\SectionProblemLink{sec:problem19-smart-manufacturing-contract-c1-c3}{Smart Manufacturing Contract (C1--C3)} \\
		\SectionProblemLink{sec:problem20-distribution-centre-automation-c1-c3}{Distribution Centre Automation (C1--C3)} \\
		\SectionProblemLink{sec:problem21-electric-vehicle-charging-network-c1-c3}{Electric Vehicle Charging Network (C1--C3)}
	}

	\subsection{Smart Manufacturing Contract (C1--C3)}
	\label{sec:problem19-smart-manufacturing-contract-c1-c3}

	A manufacturer must select conventional or automated assembly before conditions become strong demand with favourable inputs or limited demand with expensive inputs. Forecast production is \(8\) batches for conventional assembly in the favourable state and \(5\) in the expensive-input state, versus \(11\) and \(6\) batches for automation.

	Conventional batches earn a system-based \(12\) thousand USD each and automated batches \(15\); market conditions add \(4\) per batch in the favourable state or \(2\) in the expensive-input state. The contract also pays fixed revenue of \(18\) thousand USD for conventional assembly or \(28\) for automation, while a market incentive adds fixed revenue of \(14\) or \(7\) in the respective states.

	Operating rates are \(5\) thousand USD per conventional batch and \(6\) per automated batch. Materials add \(2\) per batch under favourable inputs or \(4\) under expensive inputs. Setup costs are fixed at \(24\) thousand USD for conventional assembly and \(38\) for automation, while compliance adds a fixed \(9\) in the favourable state or \(15\) in the expensive-input state.

	\textbf{Questions/Tasks.}
	\begin{enumerate}
		\CriterionTask
		{subsec:problem19-smart-manufacturing-contract-c1}{1}
		{Interpreting decision alternatives, events, consequences, and states.}

		\CriterionTask
		{subsec:problem19-smart-manufacturing-contract-c3}{3}
		{Calculating profit and building the payoff table.}
	\end{enumerate}

	\subsubsection*{C1 - Interpreting decision alternatives, events, consequences, and states.}
	\phantomsection
	\label{subsec:problem19-smart-manufacturing-contract-c1}

	\DecisionElementsTable
	{Use the conventional assembly system; Use the automated assembly system}
	{Strong demand with favourable input conditions; Limited demand with expensive inputs}
	{Demand and input-market conditions revealed after the production-system decision}
	{Profit in thousand USD for each production-system--state pair}

	Revenue contains alternative-based and state-based variable components and
	alternative-based and state-based fixed components:
	\[
	R(A_i,S_j)
	=
	q(A_i,S_j)
	\left[
	r(A_i)+r(S_j)
	\right]
	+
	f_r(A_i)
	+
	f_r(S_j).
	\]

	Cost also contains alternative-based and state-based variable components and
	alternative-based and state-based fixed components:
	\[
	C(A_i,S_j)
	=
	q(A_i,S_j)
	\left[
	c(A_i)+c(S_j)
	\right]
	+
	f_c(A_i)
	+
	f_c(S_j).
	\]

	Therefore,
	\[
	\begin{aligned}
		P(A_i,S_j)
		={}&
		\left\{
		q(A_i,S_j)
		\left[
		r(A_i)+r(S_j)
		\right]
		+
		f_r(A_i)
		+
		f_r(S_j)
		\right\}
		\\
		&-
		\left\{
		q(A_i,S_j)
		\left[
		c(A_i)+c(S_j)
		\right]
		+
		f_c(A_i)
		+
		f_c(S_j)
		\right\}.
	\end{aligned}
	\]

	\secInicio[sec:problem19-smart-manufacturing-contract-c1-c3][sec:complete-four-component-models]

	\subsubsection*{C3 - Calculating profit and building the payoff table.}
	\phantomsection
	\label{subsec:problem19-smart-manufacturing-contract-c3}

	For every alternative--state pair, calculate
	\[
	\begin{aligned}
		R(A_i,S_j)
		&=
		q(A_i,S_j)
		\left[
		r(A_i)+r(S_j)
		\right]
		+
		f_r(A_i)
		+
		f_r(S_j),
		\\
		C(A_i,S_j)
		&=
		q(A_i,S_j)
		\left[
		c(A_i)+c(S_j)
		\right]
		+
		f_c(A_i)
		+
		f_c(S_j),
		\\
		P(A_i,S_j)
		&=
		R(A_i,S_j)-C(A_i,S_j).
	\end{aligned}
	\]

	\begin{center}
		\small
		\begin{tabular}{llll}
			\toprule
			Cell
			& \(R(A_i,S_j)\)
			& \(C(A_i,S_j)\)
			& \(P(A_i,S_j)=R(A_i,S_j)-C(A_i,S_j)\) \\
			\midrule
			\((A_1,S_1)\)
			& \(8(12+4)+18+14=160\)
			& \(8(5+2)+24+9=89\)
			& \(160-89=71\) \\
			
			\((A_1,S_2)\)
			& \(5(12+2)+18+7=95\)
			& \(5(5+4)+24+15=84\)
			& \(95-84=11\) \\
			
			\((A_2,S_1)\)
			& \(11(15+4)+28+14=251\)
			& \(11(6+2)+38+9=135\)
			& \(251-135=116\) \\
			
			\((A_2,S_2)\)
			& \(6(15+2)+28+7=137\)
			& \(6(6+4)+38+15=113\)
			& \(137-113=24\) \\
			\bottomrule
		\end{tabular}
	\end{center}

	The resulting payoff table is:

	\begin{center}
		\begin{tabular}{lrr}
			\toprule
			Alternative
			& Strong demand, favourable inputs
			& Limited demand, expensive inputs \\
			\midrule
			Conventional assembly system & 71 & 11 \\
			Automated assembly system & 116 & 24 \\
			\bottomrule
		\end{tabular}
	\end{center}

	\secInicio[sec:problem19-smart-manufacturing-contract-c1-c3][sec:complete-four-component-models]

	\subsection{Distribution Centre Automation (C1--C3)}
	\label{sec:problem20-distribution-centre-automation-c1-c3}

	A retailer must choose manual, semi-automated, or robotic distribution before order volume and energy prices resolve into high volume with favourable energy or moderate volume with expensive energy. Forecast quantities, in hundreds of order batches, are \(8\) and \(5\) for manual distribution, \(11\) and \(7\) for semi-automation, and \(14\) and \(8\) for robotics across those respective states.

	The configurations earn \(12\), \(15\), and \(18\) thousand USD per hundred order batches, respectively, because their service capabilities differ. The market adds \(4\) per hundred batches under favourable energy or \(2\) under expensive energy. Service contracts provide fixed revenue of \(16\), \(24\), and \(35\) thousand USD for manual, semi-automated, and robotic distribution, respectively; delivery incentives add fixed revenue of \(15\) in the high-volume state or \(8\) in the moderate-volume state.

	Operating rates are \(5\), \(6\), and \(8\) thousand USD per hundred batches for the three respective configurations. Energy adds \(2\) per hundred batches in the favourable state or \(4\) in the expensive state. Installation creates fixed costs of \(22\), \(34\), and \(50\) thousand USD for manual, semi-automated, and robotic distribution, respectively, while compliance adds fixed costs of \(9\) under high volume or \(16\) under moderate volume.

	\textbf{Questions/Tasks.}
	\begin{enumerate}
		\CriterionTask
		{subsec:problem20-distribution-centre-automation-c1}{1}
		{Interpreting decision alternatives, events, consequences, and states.}

		\CriterionTask
		{subsec:problem20-distribution-centre-automation-c3}{3}
		{Calculating profit and building the payoff table.}
	\end{enumerate}

	\subsubsection*{C1 - Interpreting decision alternatives, events, consequences, and states.}
	\phantomsection
	\label{subsec:problem20-distribution-centre-automation-c1}

	\DecisionElementsTable
	{Use the manual distribution configuration; Use the semi-automated distribution configuration; Use the robotic distribution configuration}
	{High order volume with favourable energy prices; Moderate order volume with expensive energy}
	{Order volume and energy-market conditions revealed after the distribution-configuration decision}
	{Profit in thousand USD for each distribution-configuration--state pair}

	Revenue contains alternative-based and state-based variable components and
	alternative-based and state-based fixed components:
	\[
	R(A_i,S_j)
	=
	q(A_i,S_j)
	\left[
	r(A_i)+r(S_j)
	\right]
	+
	f_r(A_i)
	+
	f_r(S_j).
	\]

	Cost also contains alternative-based and state-based variable components and
	alternative-based and state-based fixed components:
	\[
	C(A_i,S_j)
	=
	q(A_i,S_j)
	\left[
	c(A_i)+c(S_j)
	\right]
	+
	f_c(A_i)
	+
	f_c(S_j).
	\]

	Therefore,
	\[
	\begin{aligned}
		P(A_i,S_j)
		={}&
		\left\{
		q(A_i,S_j)
		\left[
		r(A_i)+r(S_j)
		\right]
		+
		f_r(A_i)
		+
		f_r(S_j)
		\right\}
		\\
		&-
		\left\{
		q(A_i,S_j)
		\left[
		c(A_i)+c(S_j)
		\right]
		+
		f_c(A_i)
		+
		f_c(S_j)
		\right\}.
	\end{aligned}
	\]

	\secInicio[sec:problem20-distribution-centre-automation-c1-c3][sec:complete-four-component-models]

	\subsubsection*{C3 - Calculating profit and building the payoff table.}
	\phantomsection
	\label{subsec:problem20-distribution-centre-automation-c3}

	For every alternative--state pair, calculate
	\[
	\begin{aligned}
		R(A_i,S_j)
		&=
		q(A_i,S_j)
		\left[
		r(A_i)+r(S_j)
		\right]
		+
		f_r(A_i)
		+
		f_r(S_j),
		\\
		C(A_i,S_j)
		&=
		q(A_i,S_j)
		\left[
		c(A_i)+c(S_j)
		\right]
		+
		f_c(A_i)
		+
		f_c(S_j),
		\\
		P(A_i,S_j)
		&=
		R(A_i,S_j)-C(A_i,S_j).
	\end{aligned}
	\]

	\begin{center}
		\small
		\begin{tabular}{llll}
			\toprule
			Cell
			& \(R(A_i,S_j)\)
			& \(C(A_i,S_j)\)
			& \(P(A_i,S_j)=R(A_i,S_j)-C(A_i,S_j)\) \\
			\midrule
			\((A_1,S_1)\)
			& \(8(12+4)+16+15=159\)
			& \(8(5+2)+22+9=87\)
			& \(159-87=72\) \\
			
			\((A_1,S_2)\)
			& \(5(12+2)+16+8=94\)
			& \(5(5+4)+22+16=83\)
			& \(94-83=11\) \\
			
			\((A_2,S_1)\)
			& \(11(15+4)+24+15=248\)
			& \(11(6+2)+34+9=131\)
			& \(248-131=117\) \\
			
			\((A_2,S_2)\)
			& \(7(15+2)+24+8=151\)
			& \(7(6+4)+34+16=120\)
			& \(151-120=31\) \\
			
			\((A_3,S_1)\)
			& \(14(18+4)+35+15=358\)
			& \(14(8+2)+50+9=199\)
			& \(358-199=159\) \\
			
			\((A_3,S_2)\)
			& \(8(18+2)+35+8=203\)
			& \(8(8+4)+50+16=162\)
			& \(203-162=41\) \\
			\bottomrule
		\end{tabular}
	\end{center}

	The resulting payoff table is:

	\begin{center}
		\begin{tabular}{lrr}
			\toprule
			Alternative
			& High order volume
			& Moderate order volume \\
			\midrule
			Manual distribution configuration & 72 & 11 \\
			Semi-automated distribution configuration & 117 & 31 \\
			Robotic distribution configuration & 159 & 41 \\
			\bottomrule
		\end{tabular}
	\end{center}

	\secInicio[sec:problem20-distribution-centre-automation-c1-c3][sec:complete-four-component-models]

	\subsection{Electric Vehicle Charging Network (C1--C3)}
	\label{sec:problem21-electric-vehicle-charging-network-c1-c3}

	An energy company must build a neighbourhood or rapid-charging network before utilisation and electricity prices become high and favourable, moderate and standard, or low and expensive. Forecast quantities, in hundreds of charging-session batches, are \(10\), \(7\), and \(4\) for the neighbourhood network across those states, and \(14\), \(9\), and \(5\) for the rapid network.

	Neighbourhood service earns a network-based \(11\) thousand USD per hundred charging-session batches and rapid charging earns \(16\). Utilisation conditions add \(4\), \(3\), and \(2\) thousand USD per hundred batches in the high, moderate, and low states, respectively. Network service agreements provide fixed revenue of \(18\) thousand USD for the neighbourhood network or \(30\) for rapid charging, while utilisation incentives add fixed revenue of \(16\), \(10\), and \(5\) in the respective states.

	Equipment costs \(4\) thousand USD per hundred batches for the neighbourhood network or \(7\) for rapid charging. Electricity adds \(2\), \(3\), and \(5\) thousand USD per hundred batches in the respective states. Infrastructure creates a fixed cost of \(25\) thousand USD for the neighbourhood design or \(45\) for rapid charging, while state-dependent maintenance adds fixed costs of \(8\), \(12\), and \(18\), respectively. The company must preserve all eight components when computing each network--state profit.

	\textbf{Questions/Tasks.}
	\begin{enumerate}
		\CriterionTask
		{subsec:problem21-electric-vehicle-charging-network-c1}{1}
		{Interpreting decision alternatives, events, consequences, and states.}

		\CriterionTask
		{subsec:problem21-electric-vehicle-charging-network-c3}{3}
		{Calculating profit and building the payoff table.}
	\end{enumerate}

	\subsubsection*{C1 - Interpreting decision alternatives, events, consequences, and states.}
	\phantomsection
	\label{subsec:problem21-electric-vehicle-charging-network-c1}

	\DecisionElementsTable
	{Build the neighbourhood charging network; Build the rapid-charging network}
	{High utilisation with favourable electricity prices; Moderate utilisation with standard electricity prices; Low utilisation with expensive electricity}
	{Charging demand and electricity-market conditions revealed after the network-design decision}
	{Profit in thousand USD for each charging-network--state pair}

	Revenue contains alternative-based and state-based variable components and
	alternative-based and state-based fixed components:
	\[
	R(A_i,S_j)
	=
	q(A_i,S_j)
	\left[
	r(A_i)+r(S_j)
	\right]
	+
	f_r(A_i)
	+
	f_r(S_j).
	\]

	Cost also contains alternative-based and state-based variable components and
	alternative-based and state-based fixed components:
	\[
	C(A_i,S_j)
	=
	q(A_i,S_j)
	\left[
	c(A_i)+c(S_j)
	\right]
	+
	f_c(A_i)
	+
	f_c(S_j).
	\]

	Therefore,
	\[
	\begin{aligned}
		P(A_i,S_j)
		={}&
		\left\{
		q(A_i,S_j)
		\left[
		r(A_i)+r(S_j)
		\right]
		+
		f_r(A_i)
		+
		f_r(S_j)
		\right\}
		\\
		&-
		\left\{
		q(A_i,S_j)
		\left[
		c(A_i)+c(S_j)
		\right]
		+
		f_c(A_i)
		+
		f_c(S_j)
		\right\}.
	\end{aligned}
	\]

	\secInicio[sec:problem21-electric-vehicle-charging-network-c1-c3][sec:complete-four-component-models]

	\subsubsection*{C3 - Calculating profit and building the payoff table.}
	\phantomsection
	\label{subsec:problem21-electric-vehicle-charging-network-c3}

	For every alternative--state pair, calculate
	\[
	\begin{aligned}
		R(A_i,S_j)
		&=
		q(A_i,S_j)
		\left[
		r(A_i)+r(S_j)
		\right]
		+
		f_r(A_i)
		+
		f_r(S_j),
		\\
		C(A_i,S_j)
		&=
		q(A_i,S_j)
		\left[
		c(A_i)+c(S_j)
		\right]
		+
		f_c(A_i)
		+
		f_c(S_j),
		\\
		P(A_i,S_j)
		&=
		R(A_i,S_j)-C(A_i,S_j).
	\end{aligned}
	\]

	\begin{center}
		\small
		\begin{tabular}{llll}
			\toprule
			Cell
			& \(R(A_i,S_j)\)
			& \(C(A_i,S_j)\)
			& \(P(A_i,S_j)=R(A_i,S_j)-C(A_i,S_j)\) \\
			\midrule
			\((A_1,S_1)\)
			& \(10(11+4)+18+16=184\)
			& \(10(4+2)+25+8=93\)
			& \(184-93=91\) \\
			
			\((A_1,S_2)\)
			& \(7(11+3)+18+10=126\)
			& \(7(4+3)+25+12=86\)
			& \(126-86=40\) \\
			
			\((A_1,S_3)\)
			& \(4(11+2)+18+5=75\)
			& \(4(4+5)+25+18=79\)
			& \(75-79=-4\) \\
			
			\((A_2,S_1)\)
			& \(14(16+4)+30+16=326\)
			& \(14(7+2)+45+8=179\)
			& \(326-179=147\) \\
			
			\((A_2,S_2)\)
			& \(9(16+3)+30+10=211\)
			& \(9(7+3)+45+12=147\)
			& \(211-147=64\) \\
			
			\((A_2,S_3)\)
			& \(5(16+2)+30+5=125\)
			& \(5(7+5)+45+18=123\)
			& \(125-123=2\) \\
			\bottomrule
		\end{tabular}
	\end{center}

	The resulting payoff table is:

	\begin{center}
		\begin{tabular}{lrrr}
			\toprule
			Alternative
			& High utilisation
			& Moderate utilisation
			& Low utilisation \\
			\midrule
			Neighbourhood charging network & 91 & 40 & \(-4\) \\
			Rapid-charging network & 147 & 64 & 2 \\
			\bottomrule
		\end{tabular}
	\end{center}

	The negative payoff for the neighbourhood charging network under low
	utilisation represents a loss because the revenue earned does not cover the
	variable and fixed costs.

	\secInicio[sec:problem21-electric-vehicle-charging-network-c1-c3][sec:complete-four-component-models]

\end{document}
